\section{Model and transfer relations}
\label{sec:model}

The relations in this section are classical
\cite{small1972,small1972direct} and are restated only to fix notation and
the peak/rms conventions used by the gates.

With terminal voltage $e(t)$, current $i(t)$, cone displacement $x(t)$,
electrical side $e=\Re i+\Bl\dot x$ (A2), force $F=\Bl i$, and total
stiffness $k_t=1/\Cms+\rho_0c^2\Sd^2/\Vb=(1+\alpha)/\Cms$ with
$\alpha\equiv\Vas/\Vb$:

\begin{equation}
	\boxed{\
		\ddot{x}+\sigma\dot{x}+\wc^2 x = \frac{\Bl}{\Re \Mms}\,e(t)\ }
	\label{eq:ode}
\end{equation}

where $\sigma\equiv\sm+\se$, $\sm\equiv\Rms/\Mms$,
$\se\equiv\Bl^2/(\Re\Mms)$, and $\wc^2\equiv k_t/\Mms$. Free air:
$\ws=2\pi F_s=\sqrt{1/(\Mms\Cms)}$, $\Qes=\ws/\se$, $\Qms=\ws/\sm$,
$\Qts=\ws/\sigma$. In box: $\wc=\ws\sqrt{1+\alpha}$, $\Qtc=\wc/\sigma$.

\subsection{Alignment line}
\label{sec:model-alignment}

Because $\sigma$ contains no stiffness term, it is enclosure-independent,
which gives the sealed alignment relation \cite{small1972}

\begin{equation}
	\boxed{\
		\frac{\wc}{\Qtc} = \frac{\ws}{\Qts} = \sigma
			\quad\Longleftrightarrow\quad
		F_c = \Qtc\,\frac{F_s}{\Qts}\ }
	\label{eq:alignment}
\end{equation}

with $\Vb=\Vas/((\Qtc/\Qts)^2-1)$ and $\Qtc>\Qts$. A sealed driver
therefore has exactly one free alignment parameter, its position along the
line $\wc=\Qtc\sigma$, and the \emph{corner rate}

\begin{equation}
	\frac{\sigma}{2\pi}=\frac{F_s}{\Qts}
		=EBP\left(1+\frac{\Qes}{\Qms}\right)
	\label{eq:cornerrate}
\end{equation}

is the datasheet quantity that decides where the corner can be placed for a
chosen $\Qtc$. The exact expression, not the $EBP$ approximation, is used
in every gate. The envelope decay rate of every sealed alignment of a given
driver is $\sigma/2=\pi F_s/\Qts$, so the enclosure selects only the damped
frequency and the overshoot pattern, not the decay rate.

\subsection{Steady-sine relations}
\label{sec:model-steady}

For steady sine excitation with peak displacement $\hat x$, peak current
$\hat\imath$, and peak terminal voltage $\hat e$ (rms values are
$1/\sqrt2$ of these):

\begin{align}
	\hat{x}(\omega)&=
		\frac{\Bl}{\Re\Mms}\,
		\frac{\hat{e}}{\sqrt{(\wc^2-\omega^2)^2+\sigma^2\omega^2}},
		\label{eq:xofe}\\
	\hat{\imath}(\omega)&=
		\frac{\Mms}{\Bl}\,
		\hat{x}\,\sqrt{(\wc^2-\omega^2)^2+\sm^2\omega^2},
		\label{eq:iofx}\\
	\hat{p}(\omega)&=
		\frac{\rho_0\Sd\,\omega^2\hat{x}}{2\pi r}\qquad\text{(A3)} .
		\label{eq:pofx}
\end{align}

Only mechanical damping $\sm$ appears in \cref{eq:iofx}: electrical damping
is a consequence of the current, not an additional load on it. The rms
pressure that enters every SPL below is $\hat p/\sqrt2$.
